\chapter{Discussion and Evaluation}
\label{ch:discussion-evaluation}

\section{Experimental Setup and Evaluation Methodology}

To evaluate the performance, resilience, and usability of the LogiShop platform, this chapter presents a structured empirical study. The aim is to assess how well the system addresses the research gaps identified in Chapter 2, particularly geopolitical and structural route disruptions, multi-carrier data normalization, and cost-impact calculation.

The evaluation combines quantitative performance benchmarks with qualitative feature tests. For the quantitative part, the route calculation and API synchronization pipelines were tested under simulated request loads to measure response times and resource overhead. For the qualitative part, features were compared against older industry tools. Operational parameters, base contract prices, and delivery route data were drawn from anonymized telemetry provided by VN Express International Delivery Company Limited (VNEXID), which connected the algorithm implementation to real commercial constraints.

\section{Feature Evaluation and Comparative Analysis}

One contribution of this thesis is addressing the functional fragmentation in small and medium-sized enterprise (SME) logistics. Figure \ref{fig:ups-landing} and Figure \ref{fig:ups-tracking-api} show how isolated proprietary tools handle basic tracking but struggle when comparing multi-carrier alternatives or assessing spatial risk.

To position LogiShop against current logistics software, its core components were compared with \textbf{UPS.com} (a carrier-specific system) and legacy \textbf{Enterprise Logistics Management Systems (LMS)} used by large freight aggregators.

\begin{table}[H]
\centering
\caption{Feature Comparison Matrix: LogiShop vs. UPS.com}
\label{tab:feature-comparison-short}
\newcolumntype{Y}{>{\raggedright\arraybackslash}X} 
\begin{tabularx}{\textwidth}{p{4.5cm} Y Y} % Đã tăng nhẹ độ rộng cột tiêu chí và chỉnh còn 2 cột thông minh Y Y
\toprule
\textbf{Evaluation Criteria} & \textbf{UPS.com} & \textbf{LogiShop (Proposed)} \\
\midrule
\textbf{Geospatial Map Visualization} & 
Low \newline \small (Isolated text-logs / simple marker pins) & 
\textbf{High} \newline \small (Dynamic route-density layout on map) \\
\addlinespace

\textbf{Automated Alternative Route Ingestion} & 
None \newline \small (Rigid single-network constraint) & 
\textbf{Automated} \newline \small (Dynamic graph re-routing engine) \\
\addlinespace

\textbf{Disruption Cost-Impact Analysis} & 
None \newline \small (Fixed static contract rates) & 
\textbf{Dynamic} \newline \small (Integrated Storage \& Delay fine logic) \\
\addlinespace

\textbf{Multi-Carrier Integration Capability} & 
None \newline \small (Exclusive to proprietary network) & 
\textbf{High} \newline \small (Standardized third-party RESTful APIs) \\
\addlinespace

\textbf{SME Financial Accessibility Layer} & 
High \newline \small (Free basic tracking portal) & 
\textbf{High} \newline \small (Cloud-native lightweight framework) \\
\bottomrule
\end{tabularx}
\end{table}

As Table \ref{tab:feature-comparison-short} shows, carrier systems like UPS.com are reliable for tracking within their own network, but they act as data silos. They give no visibility across carriers when a company moves cargo through several networks at once. Legacy LMS platforms can ingest multi-carrier data, but they carry high operational friction: they need expensive Electronic Data Interchange (EDI) modules and manual configuration changes to reroute paths when a flight is disrupted.

LogiShop closes this gap with a low-cost, easy-to-use option for Vietnamese SMEs. By replacing rigid single-route pathfinding with a web platform that synchronizes carrier data and applies spatial risk penalties automatically, we show that small logistics companies can build more flexible supply chains without large software budgets.

\section{Algorithmic Evaluation and Disruption Handling}

The engine behind LogiShop's resilience is the dynamic graph-reweighting method defined in Algorithm \ref{alg:disruption-weighting}. To test it, we ran a simulated geopolitical crisis based on trans-Eurasian airspace closures. A test route from Tan Son Nhat International Airport (SGN) to John F. Kennedy International Airport (JFK) was mapped onto the graph $G=(V,E)$.

Instead of failing or throwing an exception, the system re-traversed the graph. It discarded the blocked northern edges and produced a southern bypass through Singapore (SIN) and Dubai (DXB), avoiding the high-risk zone.

\subsection{Disruption Cost-Impact Reconciliation}

To keep the cost accounting consistent, we separate local edge-weight penalties from contract-level penalties. The \$26.53 USD computed in Table 4.7 is the operational overhead ($C_{\text{operational\_overhead}}$) added while the shipment passes through the active US-Iran disruption corridor.

This value is used only as an edge-weight penalty in the pathfinding layer. The contract penalty for late delivery, defined in Section 5.4, is a flat \$150.00 USD per day ($P_{\text{contractual}}$).

For a shipment delayed by $\Delta t = 4.2$ days in the affected sector, the total cost ($C_{\text{total\_disruption}}$) is:

\begin{equation}
C_{\text{total\_disruption}} = \int_{0}^{t} C_{\text{operational\_overhead}}(\tau) \, d\tau + \left( \lceil \Delta t_{\text{days}} \rceil \times P_{\text{contractual}} \right)
\end{equation}

Keeping the per-edge routing costs separate from the flat contract penalties lets the algorithm stay sensitive during optimization while still accounting for the full multi-day cost that an SME would actually pay.

\section{Asynchronous API Performance and Scalability}

A major implementation challenge was staying within carrier API rate limits while avoiding interface lag during data fetching. LogiShop pulls telemetry from several carriers at once, so the backend has to handle these requests carefully.

To find the system's limits, we stress-tested the telemetry ingestion pipelines by sending continuous webhook pulses that simulated concurrent data streams. Without throttling, these requests quickly exceeded the limits of the carrier sandbox environments and triggered HTTP 429 (Too Many Requests) errors. This confirmed that naive polling fails under realistic load: unhandled synchronous requests piled up, blocked the application event loop, and caused visible lag on the user dashboard.

When a 429 error occurs, the ingestion worker stops retrying immediately, logs the error without disrupting the dashboard, and enters a backoff loop that spaces out the following requests. This keeps LogiShop running even when external APIs throttle it, and it keeps the local order data in Supabase responsive, which helps the system scale for Vietnamese logistics companies.